% Options for packages loaded elsewhere
\PassOptionsToPackage{unicode}{hyperref}
\PassOptionsToPackage{hyphens}{url}
\PassOptionsToPackage{dvipsnames,svgnames,x11names}{xcolor}
\documentclass[
  10pt,
  fontset=fandol]{ctexart}
\usepackage{xcolor}
\usepackage[margin=16mm]{geometry}
\usepackage{amsmath,amssymb}
\setcounter{secnumdepth}{-\maxdimen} % remove section numbering
\usepackage{iftex}
\ifPDFTeX
  \usepackage[T1]{fontenc}
  \usepackage[utf8]{inputenc}
  \usepackage{textcomp} % provide euro and other symbols
\else % if luatex or xetex
  \usepackage{unicode-math} % this also loads fontspec
  \defaultfontfeatures{Scale=MatchLowercase}
  \defaultfontfeatures[\rmfamily]{Ligatures=TeX,Scale=1}
\fi
\usepackage{lmodern}
\ifPDFTeX\else
  % xetex/luatex font selection
\fi
% Use upquote if available, for straight quotes in verbatim environments
\IfFileExists{upquote.sty}{\usepackage{upquote}}{}
\IfFileExists{microtype.sty}{% use microtype if available
  \usepackage[]{microtype}
  \UseMicrotypeSet[protrusion]{basicmath} % disable protrusion for tt fonts
}{}
\makeatletter
\@ifundefined{KOMAClassName}{% if non-KOMA class
  \IfFileExists{parskip.sty}{%
    \usepackage{parskip}
  }{% else
    \setlength{\parindent}{0pt}
    \setlength{\parskip}{6pt plus 2pt minus 1pt}}
}{% if KOMA class
  \KOMAoptions{parskip=half}}
\makeatother
\setlength{\emergencystretch}{3em} % prevent overfull lines
\providecommand{\tightlist}{%
  \setlength{\itemsep}{0pt}\setlength{\parskip}{0pt}}
\usepackage{amsmath,amssymb,mathtools,needspace}
\xeCJKDeclareCharClass{CJK}{"2032,"0400->"04FF,"2160->"216B,"2460->"2473,"25A0,"25C7,"25CB,"25CF}
\IfFontExistsTF{Arial Unicode MS}{\xeCJKsetup{AutoFallBack=true}\setCJKfallbackfamilyfont{\CJKrmdefault}{Arial Unicode MS}}{}
\setcounter{secnumdepth}{0}
\setlength{\emergencystretch}{2em}
\allowdisplaybreaks
\usepackage{bookmark}
\IfFileExists{xurl.sty}{\usepackage{xurl}}{} % add URL line breaks if available
\urlstyle{same}
\hypersetup{
  pdftitle={要防止大数``吃掉''小数},
  colorlinks=true,
  linkcolor={Maroon},
  filecolor={Maroon},
  citecolor={Blue},
  urlcolor={Blue},
  pdfcreator={LaTeX via pandoc}}

\title{要防止大数``吃掉''小数}
\author{}
\date{}

\begin{document}
\maketitle

\paragraph{1.4.3 要防止大数``吃掉''小数}\label{ux8981ux9632ux6b62ux5927ux6570ux5403ux6389ux5c0fux6570}

在数值运算中, 参加运算的数有时数量级相差很大, 而计算机位数有限, 如不注意运算次序就可能出现大数``吃掉''小数的现象, 影响计算结果的可靠性.

\textbf{例1.8} 在五位十进制计算机上, 计算

\[
A = 52\ 492 + \sum_{i=1}^{1000} \delta_i,
\]

其中 \(0.1 \leqslant \delta_i \leqslant 0.9\).

\textbf{解} 把运算的数写成规格化形式, 有

\[
A = 0.524\ 92 \times 10^5 + \sum_{i=1}^{1000} \delta_i.
\]

由于在计算机上计算时要对阶, 若取 \(\delta_i = 0.9\), 对阶时 \(\delta_i = 0.000\ 009 \times 10^5\), 在五位的计算机中表示为机器数 0, 因此

\[
A = 0.524\ 92 \times 10^5 + 0.000\ 009 \times 10^5 + \cdots + 0.000\ 009 \times 10^5
\]

\[
\triangleq 0.524\ 92 \times 10^5
\] （符号 \(\triangleq\) 表示机器中相等），

结果显然不可靠, 这是由运算中出现了大数 52 492``吃掉''小数 \(\delta_i\) 所造成的. 如果计算时先把数量级相同的 1 000 个 \(\delta_i\) 相加, 最后再加 52 492, 就不会出现大数``吃''小数的现象, 这时有

\[
0.1 \times 10^3 \leqslant \sum_{i=1}^{1000} \delta_i \leqslant 0.9 \times 10^3,
\]

\[
0.001 \times 10^5 + 0.524\ 92 \times 10^5 \leqslant A \leqslant 0.009 \times 10^5 + 0.524\ 92 \times 10^5,
\]

\[
52\ 592 \leqslant A \leqslant 53\ 392.
\]

\begin{center}\rule{0.5\linewidth}{0.5pt}\end{center}

\subsection{相关原页校注（agent 补充）}\label{ux76f8ux5173ux539fux9875ux6821ux6ce8agent-ux8865ux5145}

以下校注来自本节覆盖的原页，可能也涉及同页相邻小节；原文照录与校正说明分开保留。

\subsubsection{原书 PDF 页 24 的校注}\label{ux539fux4e66-pdf-ux9875-24-ux7684ux6821ux6ce8}

\href{../pages/pdf-024.md}{完整原页转录} · \href{../../source-images/pdf-024.jpeg}{原页图像}

\begin{quote}
校注（agent 补充）：本页末尾乘法次数公式的说明续下页``次乘法和 \(n\) 次加法''。
\end{quote}

\end{document}
